\chapter{Integral Calculus}

\section{Area, Distance, and Integral Approximation}

\subsection{Approximating Areas under Nonnegative Curves}

Let $f : [a,b] \to \mathbb{R}$ be a nonnegative bounded function. Our goal is to compute, or more precisely to define in a rigorous way, the area of the region
\[
S = \{ (x,y) \in \mathbb{R}^2 : x \in [a,b], \ 0 \leq y \leq f(x) \}
\]
lying between the graph of $f$ and the horizontal axis. Before attempting a general definition, it is instructive to consider the concrete example $f(x) = x^2$ on the interval $[0,1]$, and to ask how the area beneath this parabola might be approximated.

The natural idea is to divide $S$ into $n$ vertical strips $S_i$ and to approximate each strip by a rectangle. Two immediate choices for the height of the approximating rectangle present themselves: one may take the height equal to the value of $f$ at the right endpoint of the strip, or equal to the value of $f$ at the left endpoint of the strip. These two choices produce two sequences of approximations, which we may denote by $R_n$ and $L_n$ respectively. As $n$ increases, both $L_n$ and $R_n$ appear, from geometric inspection, to become increasingly accurate approximations to the true area of $S$; this suggests that the area $A$ should be defined as the common limit of the sums of the areas of the approximating rectangles, as the number of strips tends to infinity.

More generally, each strip $S_i$ may be approximated by a rectangle having the same base as the strip, but whose height is given by $f(x_i^*)$, where $x_i^*$ denotes an arbitrary point chosen within the $i$-th subinterval $[x_{i-1}, x_i]$; this leads to the approximation
\[
A \approx \sum_{i=1}^{n} f(x_i^*) \, \Delta x.
\]

In the particular example under consideration, since $f$ is increasing, one has the sandwich inequality
\[
L_n \leq \sum_{i=1}^{n} f(x_i^*) \, \Delta x \leq R_n
\]
for every choice of sample points $x_i^*$, which shows that
\[
A = \lim_{n \to +\infty} \sum_{i=1}^{n} f(x_i^*) \, \Delta x
\]
regardless of the particular choice of the points $x_i^*$ within their respective subintervals.


We now apply this idea to a general region $S$ of the form introduced above. One begins by subdividing the interval $[a,b]$ into $n$ strips $S_i$ of equal width $\Delta x = (b-a)/n$, and one then approximates each strip $S_i$ by a rectangle having the same base as the strip and height equal to $f(x_i^*)$, where $x_i^*$ is an arbitrary point chosen within the $i$-th subinterval $[x_{i-1}, x_i]$; the points
\[
x_1^*, x_2^*, \ldots, x_n^*
\]
obtained in this way are called the sample points. It can be proved that the limit
\[
\lim_{n \to \infty} \sum_{i=1}^{n} f(x_i^*) \, \Delta x
\]
always exists, and it can furthermore be shown that this limit takes the same value regardless of the particular choice of sample points. On the basis of this result, one defines the area $A$ of the region $S$ to be the limit of the sums of the areas of the approximating rectangles, that is,
\[
A = \lim_{n \to \infty} \sum_{i=1}^{n} f(x_i^*) \, \Delta x.
\]

\subsection{The Distance Problem}

A closely related problem, arising from kinematics rather than geometry, is the distance problem: given the velocity of an object at every instant of a certain time period, one wishes to determine the total distance travelled by the object during that period. If the velocity remains constant throughout the interval, the distance problem is solved immediately by the elementary formula
\[
\text{distance} = \text{velocity} \times \text{time}.
\]
The situation becomes considerably more delicate when the velocity
\[
\begin{aligned}
v &= f(t) \\
&\geq 0,
\end{aligned}
\]
so that the object always moves in the positive direction, varies with time over the interval
\[
\begin{aligned}
a &\leq t \\
&\leq b.
\end{aligned}
\]

To address this more general situation, one measures the velocity on a collection of very small subintervals of $[a,b]$, of the form $[t_{i-1}, t_i]$; within each such subinterval, an approximate value for the velocity is provided by $f(t_i^*)$, for an arbitrarily chosen point
\[
t_i^* \in [t_{i-1}, t_i],
\]
so that the total distance travelled is approximated by the sum
\[
\sum_{i=1}^{n} f(t_i^*) \, \Delta t.
\]

As the velocity is measured more and more frequently, these estimates become increasingly accurate, which suggests that the exact distance $d$ travelled should be defined as the limit of this expression,
\[
d = \lim_{n \to \infty} \sum_{i=1}^{n} f(t_i^*) \, \Delta t.
\]

This limit is of precisely the same nature as the one encountered in the area problem, which motivates the introduction of a single, unified concept capable of encompassing both situations, namely the definite integral.

\section{The Definite Integral}

\subsection{Riemann Sums and the Definition of the Integral}


In order to formulate the definition of the definite integral with full rigor, several preliminary notions are required. A partition $P$ of an interval $[a,b]$ is a finite sequence of numbers of the form
\[
a = x_0 < x_1 < x_2 < \cdots < x_n = b;
\]
the interval $[x_{i-1}, x_i]$ is called the $i$-th subinterval of the partition. The mesh, or norm, of a partition is defined to be the length of its longest subinterval, that is, writing $\Delta x_i := x_i - x_{i-1}$, one sets
\[
m_P := \max\{ \Delta x_i : 1 \leq i \leq n \}.
\]

A tagged partition of an interval $[a,b]$ is a partition together with a finite sequence of sample points
\[
x_1^*, x_2^*, \ldots, x_n^*,
\]
subject to the requirement that $x_i^* \in [x_{i-1}, x_i]$ for every index $i$. Given a tagged partition, the Riemann sum of $f$ with respect to this partition, with nodes $x_0 < x_1 < \cdots < x_n$ and sample points $x_1^*, x_2^*, \ldots, x_n^*$, is the quantity
\[
\sum_{i=1}^{n} f(x_i^*) \, \Delta x_i.
\]

\begin{definition}[Integrable Function]
Let $f$ be a bounded function defined on $[a,b]$. We say that $f$ is integrable on $[a,b]$ if the limit
\[
\lim_{m_P \to 0} \sum_{i=1}^{n} f(x_i^*) \, \Delta x_i
\]
exists, is finite, and gives the same value for every tagged partition $P$ of $[a,b]$, regardless of the choice of subdivision points and sample points, provided only that the mesh $m_P$ tends to zero.
\end{definition}

\begin{definition}[Definite Integral]
Let $f$ be an integrable function on $[a,b]$. The definite integral of $f$ from $a$ to $b$ is the limit of the Riemann sums of $f$ over any possible tagged partition of $[a,b]$, as the partitions become arbitrarily fine, that is,
\[
\int_a^b f(x) \, dx := \lim_{m_P \to 0} \sum_{i=1}^{n} f(x_i^*) \, \Delta x_i.
\]

\end{definition}

It is important to state the precise meaning of the limit appearing in this definition, since the index set over which the limit is taken, namely the collection of all tagged partitions, is considerably more general than the index set $n \in \mathbb{N}$ used in the theory of sequences. Precisely, the definition asserts that for every $\varepsilon > 0$ there exists $\delta > 0$ such that, for every tagged partition $P$ with $m_P < \delta$, one has
\[
\left| \int_a^b f(x) \, dx - \sum_{i=1}^{n} f(x_i^*) \, \Delta x_i \right| < \varepsilon.
\]

Several remarks concerning notation and convention are in order. The definite integral is a real number, and in particular it does not depend on the variable $x$ used to describe the integrand; this variable is therefore called a dummy or bound variable. Although the definition of the integral implicitly assumes that $a < b$, the characterization as a limit of Riemann sums continues to make sense even when $b \leq a$; in that case, one adopts the conventions
\[
\int_b^a f(x) \, dx := -\int_a^b f(x) \, dx, \qquad \int_a^a f(x) \, dx = 0.
\]
When $f \geq 0$ on $[a,b]$, the definite integral satisfies
\[
\int_a^b f(x) \, dx \geq 0
\]
and can be interpreted geometrically as the area under the curve $y = f(x)$ between $a$ and $b$. When $f$ takes on both positive and negative values on $[a,b]$, its definite integral can instead be interpreted as a net area, that is, as a difference of areas,
\[
\int_a^b f(x) \, dx = A_1 - A_2,
\]
where $A_1$ denotes the area of the region lying above the horizontal axis and below the graph of $f$, and $A_2$ denotes the area of the region lying below the horizontal axis and above the graph of $f$. Finally, returning to the distance problem discussed earlier, if $v = f(t)$ is the velocity of an object, then the total distance travelled during a time interval
\[
\begin{aligned}
a &\leq t \\
&\leq b,
\end{aligned}
\]
regardless of whether the object reverses direction, is given by
\[
d = \int_a^b |f(t)| \, dt.
\]

Figure~\ref{fig:ch6-riemann-sums} compares left-endpoint, right-endpoint, and midpoint approximations to the same area.

\begin{figure}[!htbp]
\centering
\begin{tikzpicture}
  \begin{axis}[
    width=6.3cm,height=6cm,
    axis lines=middle, xlabel={$x$}, ylabel={},
    xmin=-0.3,xmax=4.3, ymin=-0.3,ymax=4.0,
    xtick={0,0.57,1.14,1.71,2.29,2.86,3.43,4}, xticklabels=\empty, ytick=\empty,
    ]
    \foreach \i in {0,...,6}{
      \pgfmathsetmacro{\xa}{\i*0.5714}
      \pgfmathsetmacro{\xb}{(\i+1)*0.5714}
      \pgfmathsetmacro{\h}{0.5+2.2*sin(deg(0.55*\xa))+0.3*\xa}
      \addplot[fill=figblue,opacity=0.3,draw=figblue] coordinates
        {(\xa,0) (\xa,\h) (\xb,\h) (\xb,0)} -- cycle;
    }
    \addplot[black,thick,domain=0:4,samples=150]{0.5+2.2*sin(deg(0.55*x))+0.3*x};
  \end{axis}
  \begin{axis}[
    at={(6.9cm,0)},
    width=6.3cm,height=6cm,
    axis lines=middle, xlabel={$x$}, ylabel={},
    xmin=-0.3,xmax=4.3, ymin=-0.3,ymax=4.0,
    xtick={0,0.57,1.14,1.71,2.29,2.86,3.43,4}, xticklabels=\empty, ytick=\empty,
    ]
    \foreach \i in {0,...,6}{
      \pgfmathsetmacro{\xa}{\i*0.5714}
      \pgfmathsetmacro{\xb}{(\i+1)*0.5714}
      \pgfmathsetmacro{\h}{0.5+2.2*sin(deg(0.55*\xb))+0.3*\xb}
      \addplot[fill=figred,opacity=0.3,draw=figred] coordinates
        {(\xa,0) (\xa,\h) (\xb,\h) (\xb,0)} -- cycle;
    }
    \addplot[black,thick,domain=0:4,samples=150]{0.5+2.2*sin(deg(0.55*x))+0.3*x};
  \end{axis}
  \begin{axis}[
    at={(0,-5.4cm)},
    width=6.3cm,height=6cm,
    axis lines=middle, xlabel={$x$}, ylabel={},
    xmin=-0.3,xmax=4.3, ymin=-0.3,ymax=4.0,
    xtick={0,0.57,1.14,1.71,2.29,2.86,3.43,4},
    xticklabels={$x_0$,$x_1$,$x_2$,$x_3$,$x_4$,$x_5$,$x_6$,$x_7$},
    xticklabel style={font=\tiny}, ytick=\empty,
    ]
    \foreach \i in {0,...,6}{
      \pgfmathsetmacro{\xa}{\i*0.5714}
      \pgfmathsetmacro{\xb}{(\i+1)*0.5714}
      \pgfmathsetmacro{\xm}{(\xa+\xb)/2}
      \pgfmathsetmacro{\h}{0.5+2.2*sin(deg(0.55*\xm))+0.3*\xm}
      \addplot[fill=figgreen,opacity=0.3,draw=figgreen] coordinates
        {(\xa,0) (\xa,\h) (\xb,\h) (\xb,0)} -- cycle;
    }
    \addplot[black,thick,domain=0:4,samples=150]{0.5+2.2*sin(deg(0.55*x))+0.3*x};
  \end{axis}
  \begin{axis}[
    at={(6.9cm,-5.4cm)},
    width=6.3cm,height=6cm,
    axis lines=middle, xlabel={$x$}, ylabel={},
    xmin=-0.3,xmax=4.3, ymin=-0.3,ymax=4.0, xtick=\empty, ytick=\empty,
    ]
    \addplot[black,thick,domain=0:4,samples=150]{0.5+2.2*sin(deg(0.55*x))+0.3*x};
    \addplot[fill=black,opacity=0.12,draw=none,domain=0:4,samples=150] {0.5+2.2*sin(deg(0.55*x))+0.3*x} \closedcycle;
  \end{axis}
\end{tikzpicture}
\caption[Riemann sums: left, right, midpoint]{Top left: left-endpoint sums (blue). Top right: right-endpoint sums (red). Bottom left: midpoint sums (green). All three use the same partition $x_0<\dots<x_7$. As the partition is refined, the sums converge to the area under the curve, shaded in the bottom-right panel.}
\label{fig:ch6-riemann-sums}
\end{figure}

\subsection{Which Functions Are Integrable?}

Not every bounded function is integrable in the sense just defined; the function that takes the value $0$ at every rational number and the value $1$ at every irrational number, restricted to $[0,1]$, provides a standard example of a bounded function that fails to be integrable, since the corresponding Riemann sums can be made arbitrarily close to either $0$ or $1$ depending on the choice of sample points, so that no common limiting value exists. Fortunately, the following theorem shows that the vast majority of functions encountered in practice are indeed integrable.

\begin{theorem}
If $f$ is continuous on $[a,b]$, or if $f$ is monotone on $[a,b]$, or if $f$ has only a finite number of jump discontinuities on $[a,b]$, then $f$ is integrable on $[a,b]$; that is, the definite integral
\[
\int_a^b f(x) \, dx
\]
exists.
\end{theorem}

Integrability is furthermore preserved under composition with a continuous function, as recorded in the following result.

\begin{theorem}[Composition]
If $f$ is integrable on $[a,b]$ and $\varphi$ is continuous on an interval containing $f([a,b])$, then the composition $\varphi \circ f$ is integrable on $[a,b]$.
\end{theorem}

In particular, this theorem shows that both $|f|$ and $f^2$ are integrable on $[a,b]$ whenever $f$ itself is integrable, since the absolute value function and the squaring function are continuous.

\subsection{Properties of the Definite Integral}

The definite integral obeys a number of algebraic properties, closely paralleling those established earlier for series, which greatly facilitate its computation. Let $f$ and $g$ be two integrable functions on $[a,b]$, and let $\lambda \in \mathbb{R}$ be an arbitrary constant. Then the functions $f + g$ and $\lambda \cdot f$ are also integrable on $[a,b]$, and one has the following linearity properties.
\begin{itemize}
\item \textit{Constant functions}: for the constant function equal to $\lambda$ throughout $[a,b]$, one has

\begin{align*}
\int_a^b \lambda \, dx = \lambda (b-a).
\end{align*}

\item \textit{Additivity of the integrand}: for the sum of two integrable functions,

\begin{align*}
\int_a^b [f(x) + g(x)] \, dx = \int_a^b f(x) \, dx + \int_a^b g(x) \, dx.
\end{align*}

\item \textit{Homogeneity}: for the product of an integrable function with a scalar,

\begin{align*}
\int_a^b \lambda f(x) \, dx = \lambda \int_a^b f(x) \, dx.
\end{align*}

\end{itemize}

\begin{remark}
The product $f \cdot g$ of two integrable functions on $[a,b]$ is itself an integrable function. This follows as a consequence of the linearity properties just stated, together with the algebraic identity
\[
f \cdot g = \frac{1}{2}(f+g)^2 - \frac{1}{2}(f^2 + g^2),
\]
which expresses the product as a linear combination of squares of integrable functions, each of which is integrable by the Composition Theorem.
\end{remark}

A further fundamental property describes how integrals over adjacent intervals combine.

\begin{theorem}[Additivity over Adjacent Intervals]
Let $c$ be a number such that $a < c < b$. Then $f$ is integrable on $[a,c]$ and on $[c,b]$, and
\[
\int_a^c f(x) \, dx + \int_c^b f(x) \, dx = \int_a^b f(x) \, dx.
\]

\end{theorem}

When the integrand possesses a symmetry property, the corresponding integral over a symmetric interval can often be simplified considerably, as recorded in the following theorem.

\begin{theorem}[Integrals of Symmetric Functions]
Let $f$ be integrable on $[-a,a]$. If $f$ is even, then
\[
\int_{-a}^{a} f(x) \, dx = 2 \int_0^a f(x) \, dx.
\]
If $f$ is odd, then
\[
\int_{-a}^{a} f(x) \, dx = 0.
\]

\end{theorem}

Finally, a collection of comparison properties allows one to bound the value of a definite integral without computing it explicitly.

\begin{itemize}
\item \textit{Nonnegativity}: if $f(x) \geq 0$ for every $x \in [a,b]$, then

\[
\int_a^b f(x) \, dx \geq 0.
\]

\item \textit{Monotonicity}: if $f(x) \geq g(x)$ for every $x \in [a,b]$, then

\[
\int_a^b f(x) \, dx \geq \int_a^b g(x) \, dx.
\]

\item \textit{Triangle inequality}: the integral of $f$ is bounded by the integral of its absolute value,

\begin{align*}
\left| \int_a^b f(x) \, dx \right| \leq \int_a^b |f(x)| \, dx.
\end{align*}

\item \textit{Boundedness}: if

\[
k \leq f(x) \leq K
\]
throughout $[a,b]$, then
\begin{align*}
k(b-a) \leq \int_a^b f(x) \, dx \leq K(b-a).
\end{align*}

\end{itemize}

Figure~\ref{fig:atlas-signed-area} distinguishes signed integration from total geometric area.

\begin{figure}[!htbp]
\centering
\begin{tikzpicture}[>=Stealth]
\begin{axis}[width=11cm,height=5.8cm,axis lines=middle,ticklabel style={font=\small},label style={font=\small},legend style={font=\scriptsize,draw=black!20,fill=white},xlabel={$x$},ylabel={$\sin x$},xmin=0,xmax=6.6,ymin=-1.3,ymax=1.3,xtick={0,3.14159,6.28318},xticklabels={$0$,$\pi$,$2\pi$},ytick={-1,1}]
\addplot[draw=none,fill=figblue!25,domain=0:pi,samples=80] {sin(deg(x))} \closedcycle;
\addplot[draw=none,fill=figred!25,domain=pi:2*pi,samples=80] {sin(deg(x))} \closedcycle;
\addplot[figblue,very thick,domain=0:2*pi,samples=160] {sin(deg(x))};
\node at (axis cs:1.57,0.45) {$+$};
\node at (axis cs:4.71,-0.45) {$-$};
\end{axis}
\end{tikzpicture}
\caption{The graph of $\sin x$ has equal positive and negative areas on $[0,2\pi]$. Its signed integral is zero, whereas the integral of its absolute value is $4$. Blue and red shading indicate the opposite signs of the contributions.}
\label{fig:atlas-signed-area}
\end{figure}

\subsection{The Integral Mean Value Theorem}

A particularly useful consequence of the comparison properties, combined with the Intermediate Value Theorem for continuous functions, is the following result, which asserts that a continuous function attains its own average value at some point of the interval.

\begin{theorem}[Integral Mean Value Theorem]
Let $f$ be continuous on $[a,b]$. Then there exists $c \in [a,b]$ such that
\[
f(c) = \frac{1}{b-a} \int_a^b f(x) \, dx.
\]
This value is called the integral mean value of $f$ on $[a,b]$.
\end{theorem}

\begin{proof}
Let $m$ and $M$ denote respectively the minimum and maximum values of $f$ on $[a,b]$, guaranteed to exist by Weierstrass' Theorem, so that
\[
\begin{aligned}
m &\leq f(x) \\
&\leq M
\end{aligned}
\]
for every $x \in [a,b]$. By the boundedness property of the definite integral, it follows that
\[
m(b-a) \leq \int_a^b f(x) \, dx \leq M(b-a),
\]
and hence, setting
\[
y := \frac{1}{b-a} \int_a^b f(x) \, dx,
\]
we obtain
\[
\begin{aligned}
m &\leq y \\
&\leq M.
\end{aligned}
\]

Applying the Intermediate Value Theorem for continuous functions on $[a,b]$, we deduce the existence of a point $c \in [a,b]$ such that $f(c) = y$, which is precisely the desired conclusion.
\end{proof}

\section{The Fundamental Theorem of Calculus}

\subsection{The Integral Function and Its Derivative}

We now come to the central result of the entire theory, which reveals a deep and unexpected connection between the operations of differentiation and integration. Let $f$ be integrable on $[a,b]$, and fix any point $x \in [a,b]$. The quantity
\[
\int_a^x f(t) \, dt
\]
is a well-defined real number depending on the choice of $x$; letting $x$ vary throughout $[a,b]$ defines a new function. For every $x \in [a,b]$, set
\[
g(x) = \int_a^x f(t) \, dt.
\]

\begin{definition}[Integral Function]
The map $x \mapsto g(x)$ defined above is called the integral function, or area function, of $f$.
\end{definition}

The first part of the Fundamental Theorem of Calculus asserts that, whenever $f$ is continuous, its integral function is not merely continuous but is in fact differentiable, and that its derivative recovers the original function $f$.

\begin{theorem}[Fundamental Theorem of Calculus, Part I]
If $f$ is continuous on $[a,b]$, then the area function
\[
g(x) = \int_a^x f(t) \, dt
\]
is continuous and differentiable on $[a,b]$, and its derivative is given by $g'(x) = f(x)$ for every $x \in [a,b]$.
\end{theorem}

\begin{proof}
Fix $x$ and $h$ such that $x+h \in [a,b]$; we assume without loss of generality that $h > 0$. By the additivity property of the definite integral, we may write
\[
\begin{aligned}
g(x+h) - g(x)
&= \int_a^{x+h} f(t) \, dt - \int_a^x f(t) \, dt \\
&= \left[ \int_a^x f(t) \, dt + \int_x^{x+h} f(t) \, dt \right]
- \int_a^x f(t) \, dt \\
&= \int_x^{x+h} f(t) \, dt.
\end{aligned}
\]

Since $f$ is continuous, we may apply the Integral Mean Value Theorem on the interval $[x, x+h]$, obtaining a point $c_h \in [x, x+h]$ such that
\[
\int_x^{x+h} f(t) \, dt = f(c_h) \cdot h.
\]

As $h \to 0$, the point $c_h$, being squeezed between $x$ and $x+h$, satisfies $c_h \to x$; since $f$ is continuous, it follows that $f(c_h) \to f(x)$. Therefore,
\[
\lim_{h \to 0} \frac{g(x+h) - g(x)}{h} = \lim_{h \to 0} \frac{1}{h} \int_x^{x+h} f(t) \, dt = \lim_{h \to 0} f(c_h) = f(x).
\]
This shows that $g$ is differentiable at $x$ and that $g'(x) = f(x)$, as claimed.
\end{proof}

Using Leibniz notation for derivatives, the conclusion of the first part of the Fundamental Theorem of Calculus may be restated compactly as follows: if $f$ is continuous on $[a,b]$, then for every $x \in (a,b)$,
\[
\frac{d}{dx} \int_a^x f(t) \, dt = f(x).
\]

This identity is of great practical importance, since it means that the entire apparatus of differential calculus can be brought to bear on the study of the area function $g$, even when $g$ itself cannot be expressed in terms of elementary functions.

A useful extension of this result concerns integral functions whose limits of integration are themselves differentiable functions of $x$, rather than $x$ itself; given a continuous function $f$ and two differentiable functions $k$ and $h$, one may seek a formula for the derivative of
\[
\int_{h(x)}^{k(x)} f(t) \, dt,
\]
which can be obtained by combining the Fundamental Theorem of Calculus with the Chain Rule and the additivity property of the integral. It is also worth noting that continuity of $f$, while sufficient, is not strictly necessary for the continuity of the area function: if $f$ is merely assumed to be integrable, but not necessarily continuous, on $[a,b]$, one can still prove that the corresponding area function $g$ remains continuous on $[a,b]$, although in general it need no longer be differentiable.

As applications of these ideas, consider the sine integral
\[
\operatorname{Si}(x) = \int_0^x (\sin t)/t \, dt,
\]
where the integrand is continuously extended at $t=0$ by the value $1$.
The Fundamental Theorem gives
\[
\operatorname{Si}'(x)=\frac{\sin x}{x}\quad(x\neq0),
\qquad \operatorname{Si}'(0)=1.
\]
For the function
\[
G(x)=\int_0^{x^4} \sin(t^2) \, dt,
\]
the upper limit is $u(x)=x^4$. Combining the Fundamental Theorem with
the Chain Rule yields
\[
G'(x)=\sin\bigl((x^4)^2\bigr)\,4x^3
=4x^3\sin(x^8).
\]
Finally, let
\[
F(x) = \int_0^x f(t) \, dt
\]
for a continuous function $f$. If $f$ is even, then the substitution
$t=-u$ gives $F(-x)=-F(x)$, so $F$ is odd. If $f$ is odd, the same
substitution gives $F(-x)=F(x)$, so $F$ is even.

\subsection{Antiderivatives}

The Fundamental Theorem of Calculus naturally leads to the notion of an antiderivative, which plays a central role in the practical evaluation of definite integrals.

\begin{definition}
We say that $F$ is an antiderivative of $f$ on an interval $I$ if $F$ is differentiable on $I$ and $F'(x) = f(x)$ for every $x \in I$.
\end{definition}

\begin{remark}
If $F$ is an antiderivative of $f$ on $I$, then for any constant $c \in \mathbb{R}$ one has
\[
\begin{aligned}
\frac{d}{dx} [F(x) + c] &= \frac{d}{dx} [F(x)] + \frac{d}{dx} c \\
&= f(x) + 0 \\
&= f(x),
\end{aligned}
\]
so that $F + c$ is also an antiderivative of $f$ on $I$. Thus, whenever $f$ admits an antiderivative, it in fact admits infinitely many, which shows that the antiderivative of a given function is never unique.
\end{remark}

The following theorem shows that, on an interval, this non-uniqueness is fully accounted for by the addition of an arbitrary constant, and no other source of ambiguity is present.

\begin{theorem}
Given any two antiderivatives $F$ and $G$ of the same function $f$ on an interval $I$, the difference $F - G$ is constant on $I$.
\end{theorem}

\begin{proof}
For every $x \in I$ we have
\[
\begin{aligned}
\frac{d}{dx} [F(x) - G(x)] &= \frac{d}{dx}[F(x)] - \frac{d}{dx}[G(x)] \\
&= f(x) - f(x) \\
&= 0.
\end{aligned}
\]

We have thus shown that the function $F - G$ has vanishing derivative throughout the interval $I$, and hence, by the Increasing/Decreasing Test applied in both directions, $F - G$ must be constant on $I$; that is, there exists a constant $C \in \mathbb{R}$ such that $G(x) - F(x) = C$ for every $x \in I$.
\end{proof}

This theorem may be rephrased as follows: if $F$ is an antiderivative of $f$ on the interval $I$, then every other antiderivative of $f$ on $I$ is of the form $G(x) = F(x) + C$ for some arbitrary constant $C$.

\begin{remark}
The conclusion of this theorem is false, in general, on a domain that is not an interval. For instance, the most general antiderivative of $f(x) = -1/x^2$, defined for $x \neq 0$, is given by
\[
F(x) =
\begin{cases}
\dfrac{1}{x} + c_1 & x < 0, \\[1mm]
\dfrac{1}{x} + c_2 & x > 0,
\end{cases}
\]
for arbitrary and, crucially, independent constants $c_1, c_2 \in \mathbb{R}$, since the domain $\mathbb{R} \setminus \{0\}$ consists of two disjoint intervals rather than a single connected interval.
\end{remark}

\subsection{Existence of Antiderivatives}

It is natural to ask whether every function admits an antiderivative, and the Fundamental Theorem of Calculus provides a complete answer for the class of continuous functions: every continuous function $f$ has an antiderivative, and one such antiderivative $F$ is given by the definite integral of $f$ with variable upper limit,
\[
F(x) = \int_a^x f(t) \, dt,
\]
precisely as established in the first part of the Fundamental Theorem of Calculus. It should be emphasized, however, that although such an antiderivative always exists for a continuous function, it cannot always be expressed in terms of elementary functions; well-known examples of continuous functions whose antiderivatives escape elementary expression include $e^{-x^2}$, $\sin(x^2)$, $(\sin x)/x$, and $1/\ln x$.

It is also possible for non-continuous functions to possess antiderivatives, as illustrated by the function
\[
F(x) =
\begin{cases}
x^2 \sin(1/x) & x \neq 0, \\
0 & x = 0.
\end{cases}
\]
Computing its derivative directly from the definition, one finds
\[
F'(x) =
\begin{cases}
2x \sin(1/x) - \cos(1/x) & x \neq 0, \\
0 & x = 0.
\end{cases}
\]

The function $f(x) := F'(x)$ obtained in this way is not continuous at $x = 0$, since $\cos(1/x)$ oscillates without limit as $x \to 0$; nonetheless, $F$ is by construction an antiderivative of $f$ on every interval containing $x = 0$, showing that continuity of $f$ is sufficient, but not necessary, for the existence of an antiderivative.

On the other hand, not every function admits an antiderivative: a function with a jump discontinuity, such as
\[
f(x) =
\begin{cases}
-1 & x < 0, \\
1 & x \geq 0,
\end{cases}
\]
cannot be the derivative of any function on an interval containing the origin. This impossibility is a direct consequence of the following remarkable theorem, which shows that derivatives, even of merely differentiable functions, always satisfy an intermediate value property analogous to that of continuous functions, even when the derivative itself fails to be continuous.

\begin{theorem}[Darboux's Theorem]
Let $F : [a,b] \to \mathbb{R}$ be a differentiable function, with derivative function $F'(x) = f(x)$ for every $x \in [a,b]$. Then $f$ has the intermediate value property, that is, if $\lambda$ is any real number lying between $f(a)$ and $f(b)$, then there exists a point $c$ between $a$ and $b$ such that $f(c) = \lambda$.
\end{theorem}

\begin{proof}
Assume that $f(a) < f(b)$, and take any $\lambda$ such that $f(a) < \lambda < f(b)$; the case $f(a) > f(b)$ is treated analogously. Define the auxiliary function $H(x) = F(x) - \lambda x$ on $[a,b]$, whose derivative is $H'(x) = f(x) - \lambda$. Since $H$ is continuous on $[a,b]$, being differentiable, it attains its absolute minimum somewhere on $[a,b]$ by Weierstrass' Theorem; we claim that this minimum is attained at an interior point $c$ of $(a,b)$, rather than at one of the two endpoints. Indeed,
\[
H'(a) = f(a) - \lambda < 0
\]
shows that $H$ is decreasing near $a$, so the minimum cannot occur at $a$, while
\[
H'(b) = f(b) - \lambda > 0
\]
shows that $H$ is increasing near $b$, so the minimum cannot occur at $b$ either. Since the minimum is attained at an interior point $c \in (a,b)$, Fermat's Theorem yields $H'(c) = 0$. But $H'(c) = f(c) - \lambda$, so $f(c) - \lambda = 0$, that is, $f(c) = \lambda$, which proves the theorem.
\end{proof}

\section{Evaluating Definite Integrals}

\subsection{The Fundamental Theorem of Calculus, Part II}

Having established the existence of antiderivatives for continuous functions, we now present the second part of the Fundamental Theorem of Calculus, which furnishes an extraordinarily efficient method for evaluating definite integrals, provided that an antiderivative of the integrand can be found explicitly.

\begin{theorem}[Fundamental Theorem of Calculus, Part II]
If $f$ is integrable on $[a,b]$ and has an antiderivative $F$ on $[a,b]$, then
\[
\int_a^b f(x) \, dx = F(b) - F(a).
\]
In particular, this conclusion holds for every continuous function $f$.
\end{theorem}

It is worth pausing to appreciate how remarkable this result is: the definite integral
\[
\int_a^b f(x) \, dx,
\]
defined by means of a rather complicated limiting procedure involving the values of $f$ at every point of $[a,b]$, can nonetheless be computed simply by evaluating a suitable function $F$ at only the two endpoints $a$ and $b$.

\begin{proof}[Proof for continuous $f$]
Define the area function of $f$ by
\[
g(x) = \int_a^x f(t) \, dt
\]

By the first part of the Fundamental Theorem of Calculus, we know that $g' = f$, that is, $g$ is itself an antiderivative of $f$. By the definition of $g$, we have
\[
g(b) - g(a) = \int_a^b f(t) \, dt - \int_a^a f(t) \, dt = \int_a^b f(t) \, dt.
\]

If $F$ is any other antiderivative of $f$ on $[a,b]$, then $g$ and $F$ differ by a constant, so there exists $C \in \mathbb{R}$ such that $F(x) = g(x) + C$ for every $x \in [a,b]$. Consequently,
\[
F(b) - F(a) = [g(b) + C] - [g(a) + C] = g(b) - g(a) = \int_a^b f(t) \, dt,
\]
which completes the proof.
\end{proof}

\subsection{Differentiation and Integration as Inverse Processes}

The two parts of the Fundamental Theorem of Calculus, taken together, reveal that differentiation and integration are, in a precise sense, inverse operations of one another. If $f$ is continuous on $[a,b]$, the first part asserts that
\[
\frac{d}{dx} \int_a^x f(t) \, dt = f(x),
\]
while the second part asserts that, for any function $F$ that is differentiable with continuous derivative on $[a,b]$,
\[
\int_a^b F'(x) \, dx = F(b) - F(a).
\]

The first identity states that if a function $f$ is integrated and the result is subsequently differentiated, one recovers the original function $f$; the second identity states that if a function $F$ is first differentiated and the result is subsequently integrated, one recovers the original function $F$, in the form $F(b) - F(a)$. Together, these two facts justify the assertion that differentiation and integration are inverse processes, in the sense that each undoes what the other does.

This close relationship motivates the introduction of a companion notion to the definite integral, applicable to an entire interval rather than to a specific pair of endpoints.

\begin{definition}[Indefinite Integral]
Given a function $f : I \to \mathbb{R}$, where $I$ is an interval, we call the indefinite integral of $f$ on $I$ the set of all antiderivatives of $f$ on $I$. This set is denoted by the symbol
\[
\int f(x) \, dx.
\]

\end{definition}

\subsection{Change of Variable in the Definite Integral}

A further practical tool for the evaluation of definite integrals is provided by the change of variable formula, which allows a given integral to be transformed into a possibly simpler one through a suitable substitution.

\begin{theorem}[Change of Variable]
Let $f$ be continuous on $[a,b]$, and let $\varphi : [\alpha, \beta] \to [a,b]$ be an invertible function from $[\alpha,\beta]$ onto $[a,b]$. Assume further that $\varphi$ is differentiable, with $\varphi'$ continuous on $[\alpha,\beta]$. Then
\[
\int_a^b f(x) \, dx = \int_\alpha^\beta f(\varphi(t)) \cdot \varphi'(t) \, dt,
\]
where $\varphi(\alpha) = a$ and $\varphi(\beta) = b$.
\end{theorem}

\begin{proof}
Let $F$ be any antiderivative of $f$ on $[a,b]$. By the Chain Rule, the composition $F \circ \varphi$ is an antiderivative of $(f \circ \varphi) \, \varphi'$ on $[\alpha,\beta]$. Applying the second part of the Fundamental Theorem of Calculus to the right-hand side, we obtain
\[
\begin{aligned}
\int_\alpha^\beta f(\varphi(t)) \, \varphi'(t) \, dt
&= (F \circ \varphi)(\beta) - (F \circ \varphi)(\alpha) \\
&= F(\varphi(\beta)) - F(\varphi(\alpha)) \\
&= F(b) - F(a).
\end{aligned}
\]

Applying the second part of the Fundamental Theorem of Calculus once more, this time to the left-hand side, we likewise obtain
\[
\int_a^b f(x) \, dx = F(b) - F(a).
\]
Since both integrals equal $F(b) - F(a)$, the two integrals coincide, which proves the theorem.
\end{proof}

\section{Antiderivatives of Elementary Functions}

\subsection{Domains, Decomposition, and Chain-Rule Integrals}

The primitives, or antiderivatives, of elementary functions are derived directly from the differentiation rules established for those functions. It is essential, however, to exercise mathematical rigor concerning the domains on which these primitives are valid. A primitive $F$ of a function $f$ on an interval $(a,b)$ must, by definition, be differentiable, and consequently continuous, at every point of that interval; standard integration formulas are therefore valid only on intervals on which the corresponding primitive is indeed well defined and differentiable.

A particularly instructive example of this subtlety is furnished by the function $f(x) = 1/x$. Although it is common to state that the primitive of $1/x$ is $\ln|x|$, this compact notation requires careful interpretation. The function $\ln|x|$ is neither defined nor differentiable at $x=0$, and consequently it cannot serve as a primitive for $1/x$ on any interval containing the origin. The function $1/x$ is, in fact, integrable only on intervals that do not include zero, that is, on intervals of the form $(a,b)$ with either $0 \leq a < b$ or $a < b \leq 0$. On an interval where $x > 0$, the correct primitive is $\ln(x)$, while on an interval where $x < 0$, the correct primitive is $\ln(-x)$; the compact notation $\ln|x|$ serves merely as a convenient shorthand for describing both cases simultaneously, but the underlying domain restrictions must always be respected when applying this formula.


The linearity of the operation of differentiation translates directly into a corresponding linearity for the operation of integration. Recall that the derivative of a linear combination of functions equals the corresponding linear combination of their derivatives, that is,
\[
\frac{d}{dx} \big[ \alpha F(x) + \beta G(x) \big] = \alpha F'(x) + \beta G'(x).
\]

Conversely, if $f(x)$ and $g(x)$ are functions admitting primitives $F(x)$ and $G(x)$ respectively, then the linear combination $\alpha f(x) + \beta g(x)$ admits the primitive $\alpha F(x) + \beta G(x)$. This observation yields the integration formula by decomposition,
\[
\int \big[ \alpha f(x) + \beta g(x) \big] \, dx = \alpha \int f(x) \, dx + \beta \int g(x) \, dx.
\]


The Chain Rule for differentiation provides a powerful tool for identifying primitives of composite functions. Let $F$ be a primitive of $f$ on an interval, and let $g$ be a differentiable function; the derivative of the composite function $F(g(x))$ is then given by
\[
\begin{aligned}
\frac{d}{dx} F(g(x)) &= F'(g(x)) \cdot g'(x) \\
&= f(g(x)) \cdot g'(x).
\end{aligned}
\]

Integrating both sides of this identity with respect to $x$ yields the substitution formula in its differential form,
\[
\int f(g(x)) \, g'(x) \, dx = F(g(x)) + C.
\]

This formula allows for the direct computation of what are commonly called almost immediate integrals, namely integrals whose integrand matches the form $f(g(x)) g'(x)$; whenever the inner function $g(x)$ and its derivative $g'(x)$ can be recognized within the integrand, the corresponding primitive can be written down immediately, without any further algebraic manipulation.

\begin{example}
Consider the integral
\[
\int x^3 (x^4+1)^5 \, dx.
\]

Setting $g(x) = x^4+1$, one has $g'(x) = 4x^3$, and the factor $x^3$ appearing in the integrand is proportional to $g'(x)$; taking $f(t) = t^5$, whose primitive is $F(t) = t^6/6$, and adjusting for the constant of proportionality, one obtains
\[
\int x^3 (x^4+1)^5 \, dx = \frac{1}{4} \int 4x^3 (x^4+1)^5 \, dx = \frac{1}{4} F(x^4+1) + C = \frac{1}{24} (x^4+1)^6 + C.
\]

\end{example}

By choosing specific elementary functions for $f(y)$ in the general formula above, one obtains a small table of frequently occurring almost immediate integrals. If $f(y) = y^\alpha$ with $\alpha \neq -1$, then
\[
F(y) = y^{\alpha+1}/(\alpha+1),
\]
yielding the power rule for almost immediate integrals,
\[
\int [g(x)]^\alpha \, g'(x) \, dx = \frac{[g(x)]^{\alpha+1}}{\alpha+1} + C.
\]
If instead $f(y) = 1/y$, then $F(y) = \ln|y|$, yielding the corresponding logarithmic rule,
\[
\int \frac{g'(x)}{g(x)} \, dx = \ln|g(x)| + C.
\]

\subsection{Integration by Substitution}

The method of integration by substitution generalizes the technique of almost immediate integrals to situations in which the integrand does not directly match the form $f(g(x)) g'(x)$, and is used whenever an integral proves difficult to evaluate in its original variable $x$. The method proceeds by introducing a new variable $t$ through a change of variables $x = g(t)$, in the hope that the resulting integral becomes simpler to evaluate. Formally, if $x = g(t)$, then $dx = g'(t) \, dt$, and the integral transforms according to
\[
\int f(x) \, dx = \int f(g(t)) \, g'(t) \, dt.
\]

Once the integral on the right-hand side has been evaluated with respect to $t$, one substitutes back $t = g^{-1}(x)$ in order to express the final result in terms of the original variable $x$.

\begin{example}
To compute
\[
\int \sqrt{2x+1} \, dx,
\]
we set $u = 2x+1$, so that $du = 2 \, dx$, that is, $dx = (1/2) \, du$. The integral becomes
\[
\int \sqrt{u} \cdot \frac{1}{2} \, du = \frac{1}{2} \int u^{1/2} \, du = \frac{1}{2} \cdot \frac{2}{3} u^{3/2} + C = \frac{1}{3} (2x+1)^{3/2} + C.
\]

\end{example}

\subsection{Integration by Parts}

Integration by parts is derived from the Leibniz rule for the differentiation of a product of two functions. Recall that
\[
\frac{d}{dx} \big[ f(x) g(x) \big] = f'(x) g(x) + f(x) g'(x).
\]
Integrating both sides of this identity with respect to $x$ yields
\[
f(x) g(x) = \int f'(x) g(x) \, dx + \int f(x) g'(x) \, dx,
\]
and rearranging the terms produces the formula for integration by parts,
\[
\int f(x) g'(x) \, dx = f(x) g(x) - \int f'(x) g(x) \, dx,
\]
often written more compactly, using differential notation, as
\[
\int u \, dv = uv - \int v \, du.
\]

\begin{example}
Consider the integral
\[
\int x \sin(x) \, dx.
\]

Setting $f(x) = x$ and $g'(x) = \sin(x)$, one has $f'(x) = 1$, and one may choose $g(x) = -\cos(x)$. Applying the formula for integration by parts,
\[
\int x \sin(x) \, dx = x(-\cos(x)) - \int 1 \cdot (-\cos(x)) \, dx = -x\cos(x) + \int \cos(x) \, dx,
\]
and evaluating the remaining elementary integral yields the final result, $-x\cos(x) + \sin(x) + C$.
\end{example}

This technique proves particularly useful for integrals involving products of polynomials with trigonometric functions, such as $\sin(mx)$ or $\cos(mx)$, or with exponential functions of the form $e^{mx}$; in some cases, the formula for integration by parts must be applied repeatedly, in succession, before the integral is finally resolved.

\subsection{Rational Functions and Partial Fractions}

Rational functions are quotients of polynomials, and the strategy required to integrate them depends critically on the relative degrees of the numerator and denominator, as well as on the algebraic structure of the denominator.


For integrals of the form
\[
\int (ax+b)/(cx+d) \, dx,
\]
one must first perform polynomial division whenever the degree of the numerator is greater than or equal to the degree of the denominator. In the present case, since both numerator and denominator have degree one, the fraction may be rewritten directly as
\[
\frac{ax+b}{cx+d} = \frac{a}{c} + \frac{b - \tfrac{ad}{c}}{cx+d}.
\]

Applying the linearity of the integral, together with the substitution $u = cx+d$, for which $du = c \, dx$, one obtains the explicit antiderivative
\[
\int \frac{ax+b}{cx+d} \, dx = \frac{a}{c} x + \frac{bc-ad}{c^2} \ln|cx+d| + C.
\]

For integrals of the form
\[
\int P(x)/(ax^2+bx+c) \, dx,
\]
where $P(x)$ is a polynomial of degree less than two, the appropriate technique depends on the sign of the discriminant $\Delta = b^2 - 4ac$ of the denominator; we now examine the two resulting cases in turn.

If the discriminant is positive, $\Delta > 0$, the denominator factors into two distinct real roots, and the appropriate technique is partial fraction decomposition. As an illustration, consider the integral
\[
\int 1/(x^2-x-2) \, dx.
\]
The roots of $x^2-x-2=0$ are $x=2$ and $x=-1$, so that the integrand decomposes as
\[
\frac{1}{(x-2)(x+1)} = \frac{A}{x-2} + \frac{B}{x+1}
\]
for suitable constants $A$ and $B$; clearing denominators gives
\[
1 = A(x+1) + B(x-2),
\]
and setting $x=2$ yields $A = 1/3$, while setting $x=-1$ yields $B = -1/3$. Integrating the resulting decomposition term by term, one finds
\[
\int \left( \frac{1/3}{x-2} - \frac{1/3}{x+1} \right) dx = \frac{1}{3} \ln|x-2| - \frac{1}{3} \ln|x+1| + C.
\]

If instead the discriminant is negative, $\Delta < 0$, the denominator cannot be factored into real linear terms, and the appropriate technique is to complete the square within the denominator. For a general integral of the form
\[
\int 1/(ax^2+bx+c) \, dx,
\]
one first manipulates the numerator, if it is of first degree, by splitting the integral into a part resembling $f'(x)/f(x)$ together with a constant remainder; for a purely constant numerator, one relies on the fundamental formula
\[
\int \frac{1}{x^2+a^2} \, dx = \frac{1}{a} \arctan\!\left( \frac{x}{a} \right) + C.
\]
As an illustration, consider the integral
\[
\int 1/(4x^2+2x+1) \, dx.
\]
Completing the square in the denominator, one finds
\[
\begin{aligned}
4x^2+2x+1 &= 4\left(x^2 + \frac{1}{2}x\right)+1 \\
&= 4\left[\left(x+\frac{1}{4}\right)^2 - \frac{1}{16}\right]+1 \\
&= 4\left(x+\frac{1}{4}\right)^2 + \frac{3}{4}.
\end{aligned}
\]
Setting $u = x+1/4$, so that $du = dx$, the integral becomes
\[
\int \frac{1}{4u^2 + 3/4} \, du = \int \frac{1}{\tfrac{3}{4}\left( \tfrac{16}{3} u^2 + 1 \right)} \, du,
\]
and the further substitution $v=4u/\sqrt{3}$ gives
\[
\begin{aligned}
\int \frac{\dd u}{4u^2+3/4}
&=\frac{1}{\sqrt{3}}\int\frac{\dd v}{1+v^2}\\
&=\frac{1}{\sqrt{3}}\arctan v+C.
\end{aligned}
\]
Substituting back first $v=4u/\sqrt{3}$ and then $u=x+1/4$
produces the explicit result
\[
\boxed{\int\frac{\dd x}{4x^2+2x+1}
=\frac{1}{\sqrt{3}}
\arctan\!\left(\frac{4x+1}{\sqrt{3}}\right)+C.}
\]


The two cases examined above extend to a general strategy applicable to an arbitrary rational function. If the degree of the numerator is greater than or equal to the degree of the denominator, one performs polynomial long division first, obtaining a polynomial part, which is integrated immediately term by term, together with a proper rational function, that is, one in which the degree of the numerator is strictly smaller than the degree of the denominator. The resulting proper rational function is then decomposed into partial fractions according to the factorization of the denominator into linear and irreducible quadratic factors, according to the general decomposition
\[
\frac{P(x)}{Q(x)} = \sum_i \frac{A_i}{(x-r_i)^{k_i}} + \sum_j \frac{B_j x + C_j}{(x^2+p_j x+q_j)^{m_j}}.
\]

Each term of this decomposition is then integrated separately, using the logarithmic rule established above for the linear factors, and the arctangent and substitution techniques established above for the irreducible quadratic factors.

\section{Improper Integrals}

\subsection{Improper Integrals of Type I}

In defining the definite integral
\[
\int_a^b f(x) \, dx,
\]
we have dealt exclusively with a bounded function $f$ on a bounded interval $[a,b]$. We now extend this concept to the case in which the interval of integration is infinite, of the form $[a, \infty)$; an integral of this kind is called an improper integral, and is denoted by
\[
\int_a^{\infty} f(x) \, dx.
\]

To motivate the appropriate definition, consider the region $S$ that lies under the graph of $y = f(x)$, with $f \geq 0$, above the horizontal axis, and to the right of $x = a$, that is,
\[
S = \{ (x,y) \in \mathbb{R}^2 : x \geq a, \ 0 \leq y \leq f(x) \}.
\]

One might initially suspect that, since $S$ extends infinitely far to the right, its area must necessarily be infinite; a concrete example shows that this intuition is, in fact, mistaken. Consider the infinite region lying under the curve $y = 1/x^2$, above the horizontal axis, and to the right of the line $x = 1$. The area of the portion of this region lying to the left of the line $x = b$ is given by
\[
A(b) = \int_1^b \frac{1}{x^2} \, dx = \left[ -\frac{1}{x} \right]_{x=1}^{x=b} = 1 - \frac{1}{b}.
\]
The limiting area is determined by
\[
\lim_{b \to \infty} A(b) = 1
\]

Thus, the area of the shaded region approaches the finite value $1$ as $b \to \infty$, and it is therefore natural to say that the area of the entire infinite region $S$ equals $1$, despite the unbounded extent of the region itself.

Guided by this example, we now formulate the general definition of an improper integral for a bounded function $f$, not necessarily positive, defined on $[a, \infty)$.

\begin{definition}[Improper Integral of Type I]
Suppose that the following integral exists for every $t > a$:
\[
\int_a^t f(x) \, dx.
\]
If the limit on the right exists, we define the improper integral by
\[
\int_a^{\infty} f(x) \, dx := \lim_{t \to \infty} \int_a^t f(x) \, dx.
\]

If the limit is finite, we say that the improper integral is convergent; if the limit is infinite, we say that it is divergent.
\end{definition}

The improper integral introduced in this definition retains a geometric interpretation as an area whenever $f$ is a positive function. Thus, if $f \geq 0$ and the integral
\[
\int_a^{\infty} f(x) \, dx
\]
is convergent, one defines the area of the region
\[
S = \{ (x,y) \in \mathbb{R}^2 : x \geq a, \ 0 \leq y \leq f(x) \}
\]
to be
\[
A(S) = \int_a^{\infty} f(x) \, dx.
\]
It is instructive to contrast the convergent example above with the closely related integral
\[
\int_1^{\infty} (1/x) \, dx,
\]
for which
\[
\int_1^R\frac{\dd x}{x}=[\ln x]_1^R=\ln R\longrightarrow+\infty.
\]
Thus it diverges, and the corresponding area is infinite.

\begin{remark}
Both $1/x^2$ and $1/x$ tend to $0$ as $x \to \infty$, yet the corresponding improper integrals exhibit opposite behavior: the former converges, while the latter diverges. This shows that the values of $1/x$ simply do not decrease fast enough, as $x \to \infty$, for the corresponding improper integral to attain a finite value.
\end{remark}

More generally, for $p\neq1$,
\[
\int_1^{\infty} x^{-p} \, dx
\]
satisfies
\[
\begin{aligned}
\int_1^R x^{-p}\,\dd x
&=\left[\frac{x^{1-p}}{1-p}\right]_1^R\\
&=\frac{R^{1-p}-1}{1-p}.
\end{aligned}
\]
If $p>1$, then $R^{1-p}\to0$ and the integral converges to
$1/(p-1)$. If $p<1$, it diverges; the case $p=1$ is the logarithmic
case above. Finally,
\[
\int_0^{\infty} e^{-x} \, dx
\]
satisfies
\[
\int_0^R e^{-x}\,\dd x=[-e^{-x}]_0^R=1-e^{-R}\longrightarrow1,
\]
so it converges to $1$.

Figure~\ref{fig:atlas-improper-tail} visualizes the remainder left when a convergent improper integral is truncated.

\begin{figure}[!htbp]
\centering
\begin{tikzpicture}[>=Stealth]
\begin{axis}[width=11cm,height=5.8cm,axis lines=middle,ticklabel style={font=\small},label style={font=\small},legend style={font=\scriptsize,draw=black!20,fill=white},xlabel={$x$},ylabel={$1/x^2$},xmin=0,xmax=8.5,ymin=0,ymax=1.15,xtick={1,3,5,7},ytick={0.5,1}]
\addplot[draw=none,fill=figred!25,domain=3:8,samples=100] {1/x^2} \closedcycle;
\addplot[figblue,very thick,domain=1:8,samples=130] {1/x^2};
\draw[dashed,figred] (axis cs:3,0) -- (axis cs:3,0.32);
\node[above,figred] at (axis cs:3,0.32) {$R=3$};
\draw[->,figred] (axis cs:6,0.09) -- (axis cs:8,0.09);
\end{axis}
\end{tikzpicture}
\caption{For the integrable function $1/x^2$ on $[1,\infty)$, truncation at $R=3$ leaves a tail with integral $1/3$. Only the portion of that tail up to $x=8$ is shaded; the arrow indicates its continuation to infinity.}
\label{fig:atlas-improper-tail}
\end{figure}

\subsection{Convergence Tests and the Integral Test}

Just as for series, it is frequently more convenient to establish the convergence or divergence of an improper integral by comparison with another, better-understood integral, rather than by direct computation. The following two comparison results serve precisely this purpose.

\begin{theorem}[Comparison Test]
Let $f, g \geq 0$ be two bounded functions defined on $[a, \infty)$, with $f(x) \geq g(x)$ for every $x$ sufficiently large.
\[
\begin{gathered}
\text{If }\int_a^{\infty} f(x) \, dx
\text{ is convergent, then }\int_a^{\infty} g(x) \, dx
\text{ is convergent.} \\
\text{If }\int_a^{\infty} g(x) \, dx
\text{ is divergent, then }\int_a^{\infty} f(x) \, dx
\text{ is divergent.}
\end{gathered}
\]

\end{theorem}

As an application of this theorem, one may show that the integral
\[
\int_0^{\infty} e^{-x^2} \, dx
\]
is convergent, by comparing the integrand with a suitable exponential function for large values of $x$; it can in fact be shown, using tools from multivariable calculus, that the exact value of this integral is
\[
\sqrt{\pi}/2,
\]
although this precise value lies beyond the scope of the comparison method itself. A refinement of the Comparison Test, analogous to the Asymptotic Comparison Test established earlier for series, replaces the pointwise inequality between the two functions with an asymptotic equivalence as $x \to \infty$.

\begin{theorem}[Asymptotic Comparison Test]
Let $f, g \geq 0$ be two bounded functions defined on $[a, \infty)$. If $f \sim g$ as $x \to \infty$, then the integral
\[
\begin{gathered}
\int_a^{\infty} f(x) \, dx
\text{ is convergent if and only if }
\int_a^{\infty} g(x) \, dx
\text{ is convergent, and} \\
\int_a^{\infty} f(x) \, dx
\text{ is divergent if and only if }
\int_a^{\infty} g(x) \, dx
\text{ is divergent.}
\end{gathered}
\]

\end{theorem}

As an application, one may show that the integral
\[
\int_2^{\infty} \ln\!\left(1 + 1/x^2\right) dx
\]
is convergent, by observing that the integrand is asymptotically equivalent to $1/x^2$ as $x \to \infty$, and invoking the convergence of
\[
\int_1^{\infty} x^{-2} \, dx
\]
established earlier.

The theory of improper integrals is intimately connected with the theory of series, since both arise from a common geometric idea, namely the approximation of an area by a sum of rectangles. This connection can be exploited in two directions.

On one hand, one may use an improper integral to bound the sum of a series from above. Consider rectangles of unit width erected above the intervals $[n, n+1]$, with height equal to $1/n^2$; the sum of the areas of these rectangles is precisely
\[
\sum_{n=1}^{\infty} 1/n^2.
\]

If the first rectangle is excluded, the total area of the remaining rectangles is smaller than the area under the curve $y = 1/x^2$ for $x \geq 1$, which is precisely the value of the integral
\[
\int_1^{\infty} 1/x^2 \, dx.
\]
This comparison yields the explicit numerical bound
\[
\sum_{n=1}^{\infty} \frac{1}{n^2} \leq 1 + \int_1^{\infty} \frac{1}{x^2} \, dx = 2.
\]

On the other hand, an analogous comparison can be used to establish the divergence of a series. In order to estimate
\[
\sum_{n=1}^{\infty} 1/\sqrt{n},
\]
one instead uses rectangles whose tops lie above the curve
\[
y = 1/\sqrt{x},
\]
choosing the height of each rectangle equal to the value of the function at the left endpoint of the corresponding interval; this comparison yields
\[
\sum_{n=1}^{\infty} \frac{1}{\sqrt{n}} \geq \int_1^{\infty} \frac{1}{\sqrt{x}} \, dx = \infty,
\]
so that the series diverges. These two examples are particular instances of a general principle, known as the Integral Test, which relates the convergence of a series directly to the convergence of an associated improper integral.

\begin{theorem}[Integral Test for Series]
Let $f$ be a continuous, positive, decreasing function on $[b, \infty)$, and set $a_n = f(n)$. Then the series
\[
\sum_n a_n
\]
is convergent if and only if the improper integral
\[
\int_b^{\infty} f(x) \, dx
\]
is convergent.
\end{theorem}

\begin{remark}
It is not necessary that $f$ be decreasing throughout its entire domain; what matters is that $f$ be ultimately decreasing, that is, decreasing for $x$ larger than some fixed number $N$. It should also be emphasized that one should not infer from the Integral Test that the sum of the series is equal to the numerical value of the corresponding integral; the theorem provides information only about the character of the series, and not about the precise value of its sum.
\end{remark}

As applications of the Integral Test, one may study the convergence or divergence of the series
\[
\begin{gathered}
\sum_{n=2}^{\infty} 1/[n (\ln n)^q] \\
\sum_{n=1}^{\infty} 1/n^p,
\end{gathered}
\]
for positive parameters $p, q > 0$, by comparing each series with the corresponding improper integral and applying the criteria for convergence of
\[
\int^{\infty} x^{-p} \, dx
\]
established earlier.

The definition of improper integral given above extends naturally to a left semi-infinite interval, and to the entire real line. If
\[
\int_t^b f(x) \, dx
\]
exists for every number $t \leq b$, then
\[
\int_{-\infty}^{b} f(x) \, dx := \lim_{t \to -\infty} \int_t^b f(x) \, dx,
\]
provided that this limit exists. If both
\[
\begin{gathered}
\int_{-\infty}^{c} f(x) \, dx \\
\int_c^{\infty} f(x) \, dx
\end{gathered}
\]
are convergent for some real number $c$, one then defines
\[
\int_{-\infty}^{\infty} f(x) \, dx := \int_{-\infty}^{c} f(x) \, dx + \int_c^{\infty} f(x) \, dx,
\]
and it can be shown that the resulting value does not depend on the particular choice of the splitting point $c$.

\subsection{Improper Integrals of Type II}

A second, geometrically distinct type of improper integral arises not from an unbounded interval of integration, but from an unbounded integrand. Suppose that $f$ is a positive, continuous function defined on a finite interval $[a,b)$, but possesses a vertical asymptote at $b$; let $S$ denote the unbounded region lying under the graph of $f$ and above the horizontal axis, between $a$ and $b$. The area of the portion of $S$ lying between $a$ and $t$, for $t < b$, is given by
\[
A(t) = \int_a^t f(x) \, dx.
\]

If it happens that $A(t)$ approaches a definite finite number $A$ as $t \to b^-$, one says that the area of the region $S$ equals $A$, and one writes
\[
\int_a^b f(x) \, dx := \lim_{t \to b^-} \int_a^t f(x) \, dx.
\]

This geometric idea can be used to define an improper integral of Type II even when $f$ is not necessarily a positive function.

\begin{definition}[Improper Integral of Type II]
If $f$ is continuous on $[a,b)$ and has a vertical asymptote at $x = b$, then
\[
\int_a^b f(x) \, dx := \lim_{t \to b^-} \int_a^t f(x) \, dx,
\]
provided that this limit exists as a finite number. The improper integral is called convergent if the corresponding limit exists and is finite, and divergent if the limit is infinite.
\end{definition}

An entirely analogous definition applies when the vertical asymptote occurs at the left endpoint rather than at the right: if $f$ is continuous on $(a,b]$ and has a vertical asymptote at $x = a$, one sets
\[
\int_a^b f(x) \, dx := \lim_{t \to a^+} \int_t^b f(x) \, dx,
\]
whenever this limit exists as a finite number. Finally, if $f$ is continuous on $[a,b] \setminus \{c\}$ and has a vertical asymptote at an interior point $c$, one defines the improper integral over the entire interval $[a,b]$ by splitting at the singularity,
\[
\int_a^b f(x) \, dx := \int_a^c f(x) \, dx + \int_c^b f(x) \, dx,
\]
provided that both improper integrals on the right-hand side are convergent.

\begin{remark}
From this point onward, whenever the symbol
\[
\int_a^b f(x) \, dx
\]
is encountered, one must first examine the behavior of $f$ on $[a,b]$ in order to determine whether the expression represents an ordinary definite integral or an improper integral, since the two require different definitions and different tests for existence.
\end{remark}

We now evaluate the standard endpoint examples directly. First,
\[
\begin{aligned}
\int_0^1\frac{\dd x}{x}
&=\lim_{t\to0^+}\int_t^1\frac{\dd x}{x}\\
&=\lim_{t\to0^+}[\ln x]_t^1\\
&=\lim_{t\to0^+}(-\ln t)=+\infty,
\end{aligned}
\]
so the integral diverges. By contrast,
\[
\begin{aligned}
\int_0^1\ln x\,\dd x
&=\lim_{t\to0^+}[x\ln x-x]_t^1\\
&=-1-\lim_{t\to0^+}(t\ln t-t)\\
&=-1,
\end{aligned}
\]
because $t\ln t\to0$ as $t\to0^+$. More generally, for $p\neq1$,
\[
\int_0^1 x^{-p} \, dx
\]
satisfies
\[
\begin{aligned}
\int_0^1x^{-p}\,\dd x
&=\lim_{t\to0^+}\left[\frac{x^{1-p}}{1-p}\right]_t^1\\
&=\lim_{t\to0^+}\frac{1-t^{1-p}}{1-p}.
\end{aligned}
\]
If $p<1$, then $t^{1-p}\to0$ and the value is $1/(1-p)$; if
$p>1$, then $t^{1-p}\to+\infty$ and the integral diverges. The case
$p=1$ is the logarithmic divergence computed above. Thus the integral
converges exactly for $p<1$ and diverges for $p\geq1$. This behavior is
opposite to that of
\[
\int_1^{\infty} x^{-p} \, dx
\]
which converges exactly for $p>1$.

\subsection{Improper Integrals and Probability}


A particularly important application of improper integrals arises in the theory of probability, where continuous random variables are described by means of density functions defined on the entire real line. Random variables of this kind occur throughout the sciences: examples include the time one waits on hold before an agent of a company answers a telephone call, the test scores obtained by an individual on an aptitude test, the heights and weights of individuals drawn from a homogeneous population, and the lifetime of a randomly chosen battery of a given type. In the case of a battery, one might ask, for instance, what the probability is that the battery lasts between $100$ and $200$ hours; if $X$ denotes the lifetime of the battery, this probability is denoted by $P(100 \leq X \leq 200)$.

Every continuous random variable $X$ possesses a probability density function $f$, meaning that the probability that $X$ lies between two values $a$ and $b$ is found by integrating $f$ from $a$ to $b$,
\[
P(a \leq X \leq b) = \int_a^b f(x) \, dx.
\]

Since probabilities are measured on a scale from $0$ to $1$, a density function must satisfy $f \geq 0$ together with the normalization condition
\[
\int_{-\infty}^{\infty} f(x) \, dx = 1,
\]
expressed precisely as an improper integral over the whole real line. If
the first absolute moment is finite,
\[
\int_{-\infty}^{\infty}|x|f(x)\,\dd x<\infty,
\]
the distribution has a finite mean, defined by
\[
\mu = \int_{-\infty}^{\infty} x f(x) \,\dd x.
\]
Not every probability density satisfies this condition; the Cauchy
distribution is a standard example for which no finite mean exists.


Many important random phenomena encountered in practice, such as test scores on aptitude tests, the heights and weights of individuals from a homogeneous population, or the annual rainfall recorded at a given location, are well modeled by a normal distribution. This means that the probability density function of the corresponding random variable $X$ belongs to the family of functions
\[
f(x) = \frac{1}{\sigma \sqrt{2\pi}} \, e^{-\frac{(x-\mu)^2}{2\sigma^2}}.
\]

The mean of this distribution equals $\mu$, a fact that can be verified directly from the definition of the mean given above; the positive constant $\sigma$ is called the standard deviation, and it measures how spread out the values of $X$ tend to be around the mean $\mu$. The fact that
\[
\int_{-\infty}^{\infty} f(x) \, dx = 1
\]
for this particular density function can be verified rigorously by methods belonging to multivariable calculus, and lies beyond the scope of the present discussion.

As an application of the normal distribution, one may consider intelligence quotient scores, which are known to be distributed normally with mean $100$ and standard deviation $15$; one may then ask what percentage of the population has an intelligence quotient score between $85$ and $115$, and what percentage has a score above $140$, both questions being answered by evaluating suitable definite or improper integrals of the corresponding density function. As a further application, one may consider a company whose average customer waiting time before a call is answered by a representative is five minutes, and ask for the probability that a call is answered during the first minute, as well as the probability that a customer waits more than five minutes to be answered, both of which again reduce to the evaluation of integrals of an appropriate probability density function.
